\documentclass[12pt]{article}
\usepackage[utf8]{inputenc}
\usepackage[T1]{fontenc}
\usepackage[margin=1in]{geometry}
\usepackage{setspace}
\usepackage{xcolor}
\usepackage[colorlinks=true, urlcolor=blue!60!black, linkcolor=blue!60!black]{hyperref}
\setlength{\parindent}{0em}
\setlength{\parskip}{0.45em}
\renewcommand{\baselinestretch}{1.0}
\begin{document}
\begin{center}
{\Large \textit{Carson v. Makin}, 596 U.S. 767 (2022)\par}
\vspace{0.3em}
{\normalsize Case Brief --- One-Page Edition\par}
\end{center}

\textbf{Memory Jogger.} A State funding private schooling cannot exclude religious schools, whether for their status or their religious use (\href{https://www.courtlistener.com/opinion/6479298/carson-v-makin/}{CourtListener}).

\textbf{Facts.} Rural Maine school units without high schools pay tuition at the public or approved private school a parent selects, but since 1981 only ``nonsectarian'' schools qualified. The Carsons and Nelsons wished to use the benefit at accredited Christian schools that integrated faith across the curriculum and used faith-based admissions and hiring. Maine had kept the restriction even after \textit{Zelman v. Simmons-Harris} (2002) upheld neutral school-choice programs.

\textbf{Procedural History.} The district court and the First Circuit upheld the exclusion as a permissible bar on religious ``use.'' The Supreme Court reversed.

\textbf{Votes.} 6--3: Roberts for the Court; Breyer (with Kagan, and Sotomayor in part) and Sotomayor dissented.

\textbf{Holding.} The ``nonsectarian'' requirement violates the Free Exercise Clause. Exclusions from generally available private-school aid trigger strict scrutiny, and stricter church-state separation than the Establishment Clause requires is not a compelling interest.

\textbf{Reasoning.} \textit{Trinity Lutheran} and \textit{Espinoza} control; parents' independent choice satisfies \textit{Zelman}. The program was no public-school substitute, and auditing religious use would breed entanglement. \textit{Locke v. Davey} is confined to clergy training. Breyer: the Court erased the ``play in the joints.'' Sotomayor: taxpayers must now fund schools that discriminate.

\textbf{Examples.} \emph{Same-Side:} an education savings account barring religious curricula. \emph{Opposite-Side:} a State that funds only public schools. \emph{Fence-Sitter:} neutral nondiscrimination conditions on all participating schools.

\textbf{Critique.} Critics say the Court discarded the status/use line when \textit{Zelman} offered narrower ground; defenders say the Clause protects religious exercise, not mere identity, and that the old line forced theological auditing.

\textbf{Key Quotation.} ``In short, the prohibition on status-based discrimination under the Free Exercise Clause is not a permission to engage in use-based discrimination.'' --- Chief Justice Roberts, 596 U.S. 767.

\end{document}
